% ═══════════════════════════════════════════════════════════════
% LH-07c: Lehrerhinweise - Querer y Poder
% Abgeleitet aus LE-07c (Verlauf, Differenzierung, Fehler)
% ═══════════════════════════════════════════════════════════════

\documentclass[11pt, a4paper]{article}

\usepackage[utf8]{inputenc}
\usepackage[T1]{fontenc}
\usepackage[ngerman]{babel}
\usepackage{helvet}
\renewcommand{\familydefault}{\sfdefault}

\usepackage[margin=2cm]{geometry}
\usepackage{enumitem}
\usepackage{tabularx}
\usepackage{booktabs}
\usepackage{longtable}

\usepackage{xcolor}
\usepackage{amssymb}
\usepackage{tcolorbox}

\usepackage{fancyhdr}
\usepackage{lastpage}

% Farben
\definecolor{spanisch}{HTML}{c41e3a}
\definecolor{grau}{HTML}{888888}
\definecolor{hellgrau}{HTML}{f5f5f5}
\definecolor{basis}{HTML}{2a7a4b}
\definecolor{standard}{HTML}{2c5aa0}
\definecolor{erweiterung}{HTML}{7b1fa2}
\definecolor{loesung}{HTML}{c62828}

\newcommand{\fachfarbe}{spanisch}

% Variablen
\newcommand{\thema}{Querer y Poder -- Modalverben beim Einkaufen}
\newcommand{\sequenz}{07}
\newcommand{\block}{c}
\newcommand{\fach}{Spanisch}
\newcommand{\klasse}{7}
\newcommand{\dauer}{90 Min. (Doppelstunde)}

% Kopf-/Fußzeile
\pagestyle{fancy}
\fancyhf{}
\renewcommand{\headrulewidth}{0.4pt}
\renewcommand{\footrulewidth}{0.4pt}
\lhead{\fach{} -- Klasse \klasse}
\chead{\textbf{LH-\sequenz\block{} (Lehrerhinweise)}}
\rhead{}
\lfoot{}
\rfoot{Seite \thepage\ von \pageref{LastPage}}

% Differenzierungsmarker
\newcommand{\B}{\textcolor{basis}{\textbf{[B]}}}
\newcommand{\Std}{\textcolor{standard}{\textbf{[S]}}}
\newcommand{\E}{\textcolor{erweiterung}{\textbf{[E]}}}
\newcommand{\lsg}[1]{\textcolor{loesung}{\textbf{#1}}}
\newcommand{\es}[1]{\textcolor{\fachfarbe}{\textbf{#1}}}

% Boxen
\newtcolorbox{kurzplan}{
    colback=hellgrau,
    colframe=\fachfarbe,
    boxrule=1pt,
    arc=0pt,
    fonttitle=\bfseries,
    title=Kurzplan
}

\newtcolorbox{differenzierung}{
    colback=white,
    colframe=grau,
    boxrule=0.5pt,
    arc=0pt,
    fonttitle=\bfseries,
    title=Differenzierung
}

\newtcolorbox{fehlerbox}{
    colback=white,
    colframe=loesung,
    boxrule=0.5pt,
    arc=0pt,
    fonttitle=\bfseries,
    title=Häufige Fehler
}

\newtcolorbox{materialbox}{
    colback=white,
    colframe=\fachfarbe,
    boxrule=0.5pt,
    arc=0pt,
    fonttitle=\bfseries,
    title=Material-Checkliste
}

\newtcolorbox{grammatikbox}{
    colback=hellgrau,
    colframe=\fachfarbe,
    boxrule=1pt,
    arc=0pt,
    fonttitle=\bfseries,
    title=Grammatik-Übersicht
}

% ═══════════════════════════════════════════════════════════════
\begin{document}

\begin{center}
    {\Large\bfseries\textcolor{\fachfarbe}{Lehrerhinweise: \thema}}\\[0.3em]
    {\normalsize LH-\sequenz\block{} | \dauer{} | Std. 5-6 von 8}
\end{center}

\vspace{10pt}

% ═══════════════════════════════════════════════════════════════
% GRAMMATIK-ÜBERSICHT
% ═══════════════════════════════════════════════════════════════

\section*{0. Grammatik-Übersicht}

\begin{grammatikbox}
\textbf{Stammvokalwechsel: e $\rightarrow$ ie (querer) und o $\rightarrow$ ue (poder)}

\vspace{0.5em}
\begin{tabular}{|l|c|c|c|c|}
\hline
\textbf{Person} & \textbf{querer} & \textbf{Wechsel} & \textbf{poder} & \textbf{Wechsel} \\
\hline
yo & quiero & \checkmark & puedo & \checkmark \\
tú & quieres & \checkmark & puedes & \checkmark \\
él/ella & quiere & \checkmark & puede & \checkmark \\
nosotros & queremos & -- & podemos & -- \\
vosotros & queréis & -- & podéis & -- \\
ellos/ellas & quieren & \checkmark & pueden & \checkmark \\
\hline
\end{tabular}

\vspace{0.5em}
\textbf{Merksatz:} ``Wir und ihr -- bleiben fit!'' (nosotros/vosotros = kein Wechsel)

\textbf{Regel:} Bei betonter Stammsilbe wechselt der Vokal. Bei nosotros/vosotros ist die Endung betont.
\end{grammatikbox}

% ═══════════════════════════════════════════════════════════════
% 1. KURZPLAN
% ═══════════════════════════════════════════════════════════════

\section*{1. Stundenverlauf (Kurzplan)}

\textbf{1. Stunde (45 Min.)}

\begin{tabularx}{\textwidth}{|c|l|X|l|}
\hline
\textbf{Zeit} & \textbf{Phase} & \textbf{Inhalt} & \textbf{Material} \\
\hline
0--5 & Einstieg & Bildimpuls Tienda, Audio Einkaufsdialog & Bild, Audio \\
\hline
5--10 & Aktivierung & Redemittel sammeln: quiero, puede... & Tafel \\
\hline
10--20 & Input 1 & Verb \textit{querer}: Stammwechsel entdecken & Tafel, AB \\
\hline
20--30 & Input 2 & Verb \textit{poder}: Analog einführen & Tafel, AB \\
\hline
30--40 & Übung 1 & Chorisches Sprechen: Konjugation & -- \\
\hline
40--45 & Sicherung & Merkkasten besprechen, Regel formulieren & AB-07c \\
\hline
\end{tabularx}

\vspace{1em}

\textbf{2. Stunde (45 Min.)}

\begin{tabularx}{\textwidth}{|c|l|X|l|}
\hline
\textbf{Zeit} & \textbf{Phase} & \textbf{Inhalt} & \textbf{Material} \\
\hline
0--5 & Aktivierung & Quiz: Formen abfragen (quiero, puedo...) & -- \\
\hline
5--15 & Vertiefung & EA: Lückentext AB-07c (Aufg. 2) & AB-07c \\
\hline
15--25 & Anwendung & PA: Rollenspiel Verkäufer/Kunde & Rollenkarten \\
\hline
25--35 & Transfer & \E{}: Weitere Stammwechselverben & Tafel \\
\hline
35--40 & Sicherung & 2-3 Paare präsentieren Dialog & -- \\
\hline
40--45 & Abschluss & Zusammenfassung, HA & -- \\
\hline
\end{tabularx}

% ═══════════════════════════════════════════════════════════════
% 2. DIFFERENZIERUNG
% ═══════════════════════════════════════════════════════════════

\section*{2. Differenzierung}

\begin{differenzierung}
\textbf{Legende:} \B{} Basis (BR) | \Std{} Standard (MR) | \E{} Erweiterung (GY)

\vspace{0.5em}
\textbf{WICHTIG:} Diese Marker sind NUR für die Lehrkraft!
\end{differenzierung}

\vspace{1em}

\begin{tabularx}{\textwidth}{|l|X|X|X|}
\hline
& \B{} \textbf{Basis} & \Std{} \textbf{Standard} & \E{} \textbf{Erweiterung} \\
\hline
\textbf{Aufg. 1} & Zuordnung mit sichtbarer Tabelle & Zuordnung selbstständig & + eigene Sätze bilden \\
\hline
\textbf{Aufg. 2} & Wortbank mit allen 8 Formen & Wortbank mit 4 Formen & Ohne Wortbank \\
\hline
\textbf{Aufg. 3} & Dialog mit Vorlage lesen & Dialog mit Stichworten & Freier Dialog \\
\hline
\textbf{Aufg. 4} & 6 Zeilen, mit Hilfe & 8 Zeilen, selbstständig & + Stammwechsel erklären \\
\hline
\end{tabularx}

\vspace{1em}

\textbf{Konkrete Hinweise:}
\begin{itemize}
    \item \B{} LRS-SuS: Zeitzugabe (+10 Min.), Konjugationstabelle als Hilfe erlauben
    \item \B{} DaZ-SuS: Bedeutung von ``wollen'' und ``können'' vorab klären
    \item \Std{} Standardgruppe: Partnerarbeit mit Rollenkarten
    \item \E{} Leistungsstarke: Weitere Stammwechselverben (pensar, dormir), Transfer-Aufgabe
\end{itemize}

% ═══════════════════════════════════════════════════════════════
% 3. HÄUFIGE FEHLER
% ═══════════════════════════════════════════════════════════════

\section*{3. Häufige Fehler}

\begin{fehlerbox}
\begin{tabularx}{\textwidth}{|l|X|X|}
\hline
\textbf{Fehler} & \textbf{Beispiel} & \textbf{Reaktion} \\
\hline
Stammwechsel vergessen & *quero, *podo & Farbmarkierung: qu\textbf{\textcolor{loesung}{ie}}ro, p\textbf{\textcolor{loesung}{ue}}do \\
\hline
Falscher Wechsel & *puiedo (ie statt ue) & Unterschied visualisieren: e$\rightarrow$ie vs. o$\rightarrow$ue \\
\hline
Infinitiv vergessen & *Quiero compro & Struktur: Modalverb + INFINITIV! \\
\hline
nosotros mit Wechsel & *nosotros quieremos & Merksatz: ``Wir und ihr -- bleiben fit!'' \\
\hline
Aussprache ie/ue & [i-e] statt [je] & Vorsprechen: Es ist EIN Laut, nicht zwei! \\
\hline
\end{tabularx}
\end{fehlerbox}

% ═══════════════════════════════════════════════════════════════
% 4. MATERIAL-CHECKLISTE
% ═══════════════════════════════════════════════════════════════

\section*{4. Material-Checkliste}

\begin{materialbox}
\begin{itemize}[label=$\square$]
    \item AB-07c.pdf (24 Kopien)
    \item ML-07c.pdf (1× für Lehrkraft)
    \item Lehrwerk: Qué pasa 1, S. 92--94
    \item Bildimpuls: Tienda de ropa (Beamer/Ausdruck)
    \item Audio: Einkaufsdialog (CD/USB)
    \item Rollenkarten Vendedor/Cliente (12 Sets)
    \item Farbkreide/Marker für Stammwechsel
    \item Optional: LP-07c.html (Tablets/PCs)
    \item Optional: Kleidungsstücke/Bilder für Rollenspiel
\end{itemize}
\end{materialbox}

% ═══════════════════════════════════════════════════════════════
% 5. LÖSUNGEN
% ═══════════════════════════════════════════════════════════════

\section*{5. Lösungen AB-07c}

\textbf{Merkkasten querer:}
yo $\rightarrow$ quiero, tú $\rightarrow$ quieres, él/ella $\rightarrow$ quiere, nosotros $\rightarrow$ queremos

\textbf{Merkkasten poder:}
yo $\rightarrow$ puedo, tú $\rightarrow$ puedes, él/ella $\rightarrow$ puede, nosotros $\rightarrow$ podemos

\vspace{0.5em}

\textbf{Aufgabe 1:} (je 1P, max. 6P)
\begin{itemize}
    \item querer: yo $\rightarrow$ quiero, tú $\rightarrow$ quieres, él/ella $\rightarrow$ quiere
    \item poder: yo $\rightarrow$ puedo, tú $\rightarrow$ puedes, él/ella $\rightarrow$ puede
\end{itemize}

\textbf{Aufgabe 2:} (je 1P, max. 6P)
\begin{enumerate}[label=\alph*)]
    \item \lsg{quiero}
    \item \lsg{Puedo}
    \item \lsg{quiere}
    \item \lsg{podemos}
    \item \lsg{quieres}
    \item \lsg{puede}
\end{enumerate}

\textbf{Aufgabe 3:} (je 1,5P, max. 6P)
\begin{itemize}
    \item Quiero una camiseta.
    \item ¿Puedo probármela?
    \item ¿Cuánto cuesta?
    \item ¿Puedo pagar con tarjeta?
\end{itemize}

\textbf{Aufgabe 4:} (6P)
\begin{itemize}
    \item Begrüßung: 1P
    \item querer verwenden: 1P
    \item poder verwenden: 1P
    \item Preis/Bezahlung: 1P
    \item Länge (8+ Zeilen): 1P
    \item Sprachliche Korrektheit: 1P
    \item Bonus (Stammwechsel-Regel): +1P
\end{itemize}

% ═══════════════════════════════════════════════════════════════
% 6. HAUSAUFGABE
% ═══════════════════════════════════════════════════════════════

\section*{6. Hausaufgabe}

\begin{itemize}
    \item Lehrwerk S. 93, Übung 4 (querer/poder einsetzen)
    \item Vokabeln lernen: Konjugation, Einkaufs-Redemittel
    \item Optional: Lernpfad LP-07c.html bearbeiten
\end{itemize}

% ═══════════════════════════════════════════════════════════════
% 7. REDEMITTEL-KARTEN (Kopiervorlage)
% ═══════════════════════════════════════════════════════════════

\section*{7. Redemittel-Karten (Kopiervorlage)}

\begin{center}
\fbox{\begin{minipage}{0.45\textwidth}
\textbf{VENDEDOR (Verkäufer)}

\vspace{0.5em}
\begin{itemize}[leftmargin=*]
    \item ¡Buenos días/tardes!
    \item ¿En qué puedo ayudarle?
    \item ¿Qué quiere(s)?
    \item ¿De qué color?
    \item ¿Qué talla quiere?
    \item El probador está allí.
    \item Son [X] euros.
    \item ¿Algo más?
\end{itemize}
\end{minipage}}
\hfill
\fbox{\begin{minipage}{0.45\textwidth}
\textbf{CLIENTE (Kunde)}

\vspace{0.5em}
\begin{itemize}[leftmargin=*]
    \item Quiero una/un...
    \item ¿Puedo ver eso/a?
    \item ¿Puedo probármelo/la?
    \item Talla [S/M/L/XL]
    \item Me queda bien/mal.
    \item ¿Cuánto cuesta?
    \item ¿Puedo pagar con tarjeta?
    \item ¡Gracias! ¡Hasta luego!
\end{itemize}
\end{minipage}}
\end{center}

% ═══════════════════════════════════════════════════════════════
% 8. REFLEXIONSNOTIZEN
% ═══════════════════════════════════════════════════════════════

\section*{8. Reflexionsnotizen (nach Durchführung)}

\begin{tabularx}{\textwidth}{|l|X|}
\hline
\textbf{Was lief gut?} & \\[2em]
\hline
\textbf{Was lief nicht?} & \\[2em]
\hline
\textbf{Zeitmanagement} & \\[2em]
\hline
\textbf{Für nächstes Mal} & \\[2em]
\hline
\end{tabularx}

\end{document}
